\documentclass[11pt]{article}
\usepackage[margin=25mm]{geometry}
\usepackage{amsmath,amssymb,amsthm,mathtools,hyperref}
\newtheorem{proposition}{Proposition}
\newtheorem{lemma}[proposition]{Lemma}
\newcommand{\R}{\mathbb R}
\newcommand{\rankt}{\operatorname{rank}_{\otimes}}
\newcommand{\SOS}{\Sigma}
\allowdisplaybreaks
\title{From singular tensor certificates to fixed testing features}
\author{}
\date{23 October 2026}
\begin{document}
\maketitle

\noindent
The strict-Gram construction preserves small tensor rank even when
it changes the testing functions. We prove this fact, apply it to
the all-$H$ C1 certificate, and give explicit formulas for fixing
the testing functions on a bounded source and testing region.
The passage from these local regions to a full-set pencil remains
to be constructed.

\section{What the tensor criterion does and does not supply}

Write $T=K_c\otimes_{\R}K_t$ for two labelled real closed
extensions of $\R$, and let $B=B_c\otimes B_t$ be its bounded
tensor subring. There is no disjointness assumption.
For $z\in T$, $\rankt(z)$ is the smallest number of elementary
source--testing products summing to $z$.

Scheiderer's Theorem 3.10 converts a universal bound $d$ on the
tensor rank of individual square roots into a bound $d$ on the
largest PSD block. It does not prescribe the number of blocks.
Its proof constructs the algebra of testing functions at an
individual pair, then makes two finite selections: a testing-image
cover, combined by a disjoint union, and a support-admission cover,
combined by a fibre product. See Lemmas 3.15--3.17, printed
pp.~13--15 of
\href{https://arxiv.org/pdf/2004.04196}{arXiv:2004.04196v2}.
The dimensions and functions in those finite selections have not
been evaluated in the present public Q222 source.

Thus there are three different finite objects to extract:
a rank bound for each tensor root; a common finite testing space
on one algebraic cover; and the multiplication tables that give
the affine pencil. The first does not by itself print the other two.

\section{A rank bound which survives the strict-Gram step}

\begin{lemma}\label{lem:rank}
Let $G\in\operatorname{Sym}_r(B)$ have positive-definite real
residue. Suppose $g_1,\ldots,g_r\in K_c\otimes V$, where
$V\subset K_t$ is an $\R$-linear space of dimension at most $d$.
Then $g^{\mathsf T}Gg$ is a finite sum of squares whose individual
roots have tensor rank at most $d$.
If only $\rankt(g_i)\le R$ is known, the bound $rR$ follows.
\end{lemma}
\begin{proof}
The scalar positive-residue lemma TF1 writes every positive-residue
bounded tensor as a sum of squares of elementary tensors. This is
stronger than merely giving arbitrary tensor square roots.
The strict-Gram lemma TF2 first applies a real congruence and then
expresses the quadratic form using positive-residue scalar
coefficients times squares of real linear combinations of the
$g_i$. Those combinations still lie in $K_c\otimes V$.
Multiplication by an elementary tensor $a\otimes b$ places such a
combination in $K_c\otimes bV$, whose real dimension is at most
$d$. Expanding the scalar sums of squares therefore proves the
bound. Taking $V$ to be spanned by the testing factors in displayed
rank-$R$ expansions of all $g_i$ gives $\dim V\le rR$.
\end{proof}

This proof needs no bound on the number of elementary tensors in a
Gram coefficient. Such a bound affects the number of square roots,
but not their individual tensor rank. The exact source is
\href{https://github.com/anubhav70747/q222-finite-sdp/blob/0c1ce2e15ed067a6845464e2468559cea5274f2c/latex/chapters/03_tensor.tex#L44-L118}
{\texttt{03\_tensor.tex}, TF1--TF2}.

\subsection{An evaluated rank bound for the all-$H$ C1 theorem}

Use the domain of the public theorem
\href{https://github.com/anubhav70747/q222-finite-sdp/blob/0c1ce2e15ed067a6845464e2468559cea5274f2c/latex/certificates/c1_short.tex#L7253-L7780}
{\texttt{thm:c1-unification-bounded-z-all-H}}:
$x=k^2$ and $\sigma$ are positive source infinitesimals,
$z>0$ is bounded, $h/\sigma$ is arbitrary, $u,w$ are independent
testing infinitesimals, and the original strict gaps and
$D,\beta\ge0$ are retained. No state ratio is imposed.

For its $\lambda>0$, $WZ>0$ branch, choose the common state sign
so $W,Z>0$. Every root in its main Gram belongs to the
source-linear span of the following twelve testing functions:
\begin{equation}\label{eq:h12}
\begin{split}
 \mathcal H_+=\{&X,Y,Z,W,uZ,wZ,uW,uwW,u^2W,\\
              &u^{3/2}W,u^{1/2}W,\sqrt{uwWZ}\}.
\end{split}
\end{equation}
These are distinct generalized monomials. Substituting squares
for the positive variables makes them distinct ordinary monomials,
so their real span has dimension exactly twelve on the open chart.
Homogeneous state normalization does not change this independence.

Here is the entire root check. In the notation of that theorem,
\[
\begin{array}{c|l}
\text{root}&\text{testing factors needed}\\\hline
N_x&X,\ uZ\\
N_y&Y,\ wZ\\
R&wZ,\ uW\\
\Gamma_*&Z,\ uwW\\
A_\theta&W,\ u^2W\\
T_\theta&W,\ uW\\
d_1,\ldots,d_4&uZ,\ uwW,\ W,\ wZ,\ \sqrt{uwWZ}\\
U_*,K_*&u^{3/2}W,\ u^{1/2}W,\ uW,\ W.
\end{array}
\]
All omitted factors are pure source factors. In particular the
four circuit monomials $B_i$ are elementary tensors, and their
geometric mean has testing factor $\sqrt{uwWZ}$.
The twelve-root Gram is printed at
\href{https://github.com/anubhav70747/q222-finite-sdp/blob/0c1ce2e15ed067a6845464e2468559cea5274f2c/latex/certificates/c1_short.tex#L7550-L7691}
{lines 7550--7691}. Lemma~\ref{lem:rank} therefore bounds its
post-TF2 roots by twelve, rather than the coarser $12\cdot2$.

The retained squares do not enlarge this bound. Event and radial
excess roots have rank at most three; the circuit remainder roots
have rank at most two. For the angular budget, after expanding
each $U+1$ prefactor, the only root with more than two monomials is
\[
 \sqrt a\,U(V-Z_*)-(Z_*-UV)/\sqrt a
 =(\sqrt a+1/\sqrt a)UV-\sqrt a\,UZ_*-Z_*/\sqrt a.
\]
It has rank at most three: $a$ is pure source and $U,V,Z_*$ are
separated monomials. Multiplication by the other separated factors
preserves that rank. Opposite-state-sign additions are separated
nonnegative monomials. The zero-state-coordinate cases have the
same or smaller ranks. In the $WZ\le0$ branch the $d_i=B_i$ do
not even need the twelfth testing function.

For $\lambda\le0$, the Gram uses
$(N_x,N_y,R,\Gamma,A_\theta,V)$ with
$V=uw(u^2-x)W$. Its testing span is contained in
\[
 \{X,Y,Z,W,uZ,wZ,uW,uwW,u^2W,u^3wW\},
\]
of dimension ten. The source boundary $D=0$ does not cause a
rank loss problem: $\Gamma_*=\sqrt{4D/\sigma^2}\Gamma$ is then
simply zero, and no inverse of $D$ was used.

\begin{proposition}\label{prop:local12}
Conditional on the printed all-$H$ C1 certificate, its entire
stated local domain has a universal SOS representation with
individual root tensor rank at most twelve, over all labelled
real closed extensions, all allowed source faces and all states.
This is a local tensor-rank statement. It is not a bound for the
full quantum set, since the other source and projective-testing
domains have not been included in this calculation.
\end{proposition}

\section{Where a fixed testing dictionary is still lost}

The twelve functions in \eqref{eq:h12} span the roots before
strict-Gram factorization. The Gram itself depends on the testing
variables. It is therefore not a PSD matrix of source coefficients
for those twelve functions. Treating it as one would omit exactly
the state-normal and angular corrections being certified.

More explicitly, if $V=\operatorname{span}_{\R}\mathcal H_+$ and
TF1 supplies the elementary root $a(c)b(t)$, a pre-Gram root in
$K_c\otimes V$ is sent to $K_c\otimes bV$. This preserves the
bound twelve on tensor rank, but does not in general preserve
$V$. Different scalar coefficients may supply different $bV$.
The proposition above asserts the rank bound only; it does not
assert that all final roots lie in the twelve-function span.

The post-TF1 factors are also not automatically in one finite
testing space. In TF4 the scalar splitting chooses a real
$M>|y|$ after the bounded testing element $y$ is given. Its
testing square roots include $\sqrt{M+y}$ and $\sqrt{M-y}$.
The real congruence and the positive scalar budgets may likewise
be chosen separately at each source/test pair.

\begin{lemma}\label{lem:independence}
For distinct real $M_j>1$, the functions
$\sqrt{M_j+t}$ on $-1<t<1$ are linearly independent over $\R$.
Consequently the union of the allowable TF1 choices is not a
finite-dimensional testing space, even for a linear testing factor.
\end{lemma}
\begin{proof}
In $\R(t)^*/\R(t)^{*2}$ the classes of $M_j+t$ are independent:
a product of a nonempty subset has odd valuation at each of its
distinct zeros $t=-M_j$. The associated multiquadratic function
field therefore has an automorphism changing the sign of one
chosen square root and fixing the others. Apply that automorphism
to a supposed real linear relation and subtract it from the
original relation. Its chosen coefficient is zero. Repeat for
each index. A relation between the analytic positive branches on
an interval would also be a function-field relation, so the same
conclusion applies to the displayed functions.
\end{proof}

This is an obstruction to taking the union of the displayed
pointwise roots, not an obstruction to any possible finite lift.
A different certificate can fix $M$ on a uniformly bounded box.
Indeed, on a fixed rectangle
$|x(c)|\le\varepsilon/M$, $|y(t)|\le M$, with fixed real
$\varepsilon,M>0$, the identity
\[
 \varepsilon+x(c)y(t)=\tfrac12
 \bigl[(\varepsilon/M+x(c))(M+y(t))
       +(\varepsilon/M-x(c))(M-y(t))\bigr]
\]
has nonnegative separated factors. For this channel the only
testing radicals are the two fixed functions $\sqrt{M+y(t)}$
and $\sqrt{M-y(t)}$. Equality on either boundary is harmless.
Thus multiplication of a fixed $d$-function testing span uses
at most $2d$ functions for this channel. The unresolved step is
to supply finitely many such rectangles and prove their coverage,
not the scalar identity on any one of them.
An auxiliary algebraic cover can also change the available
functions, but then one must prove a global SOS identity on that
fixed cover; an existential choice of its new coordinates at
each source/test pair is insufficient.

There is a second issue at precisely the same interface.
The manuscript's conditions ``infinitesimal'', ``bounded unit''
and ``unbounded'' refer to types over $\R$; they are not fixed
real semialgebraic boxes. A finite list of such source cases is
not yet an evaluated finite list of real source/test rectangles
with fixed constants. Qualitative real-spectrum compactness can
select those data if its universal premise is valid. A compiler
must actually supply them.

\section{A fixed-box compilation which removes that obstruction}

The following statement gives an explicit target for completing
one local box, rather than adjoining all possible TF1 radicals.

\begin{lemma}[Fixed testing radicals]\label{lem:box}
Let $h_1,\ldots,h_d$ be fixed testing functions on $Y$.
Suppose a finite source-dependent quadratic form has been written
as a sum of scalar channels times squares of source-linear
combinations of the $h_i$. In every channel suppose
\[
 q(c,t)=q_0(c)+\sum_{j=1}^{m}a_j(c)b_j(t),\qquad
 q_0(c)\ge c_0>0,\quad |a_j(c)|\le A,
\]
where $m\ge1$, the $b_j$ are a fixed list of testing functions and
$c_0,A>0$ are fixed rational numbers. Put
$\gamma=c_0/(mA)$ and assume $|b_j(t)|\le\gamma$ on $Y$.
Then a common testing spanning list is
\[
 \{h_i\}_{i=1}^d\ \cup\
 \{h_i\sqrt{\gamma+b_j},\ h_i\sqrt{\gamma-b_j}:
                     1\le i\le d, 1\le j\le m\},
\]
with at most $d(1+2m)$ functions. All scalar multipliers are
nonnegative source factors. The formula is an identity over all
labelled real closed extensions satisfying the same finite bounds.
\end{lemma}
\begin{proof}
Exactly
\[
 q=(q_0-c_0)+\frac12\sum_{j=1}^{m}
 \bigl[(A+a_j)(\gamma+b_j)+(A-a_j)(\gamma-b_j)\bigr].
\]
All four factors in every bracket are nonnegative in their own
fields. Take their separate square roots and multiply by the
channel's source-linear combination of $h_i$. The source-only
term uses no new testing function. The proposed list follows.
\end{proof}

For the bounded-$z$ C1 family, the accompanying
\emph{Fixed testing functions at singular C1 contacts} carries out this
construction on the independent rectangles
\[
0<x,\sigma\le2^{-160},\qquad 0<z\le1,
\qquad |u|,|w|\le2^{-100},
\]
with the original source gaps and $D,\beta\ge0$. The ratio
$H=h/\sigma$ is unrestricted. In the positive-angular case, twelve
main roots and seven nonconstant testing monomials give
$12(1+2\cdot7)=180$ functions by Lemma~\ref{lem:box}.
The retained event, angular and circuit squares enlarge this list
to 220. A source-only shear treats the nonpositive-angular case
with 157 functions. Thus their union contains at most 377 functions
for the complete real-state polynomial on this rectangle.

The two constructions print every source bound, coefficient matrix,
square factor and testing function. They include both testing axes,
all real state signs and the permitted zero event weights.
The JSON files supply the labelled lists and coefficient prescriptions.
Extending these constructions across the remaining source and testing
domains is the outstanding global selection problem.

\section{The final finite matrix once a common cover is evaluated}

Suppose one has actually printed a finitely generated real algebra
$A$, a visible map $\phi:Y(\R)\to\R^8$, and finite subspaces
$U_1,\ldots,U_s\subset A$ such that
\[
 \mathcal Q_{222}\subset\phi(Y(\R)),\qquad
 \phi^*(P)\subset\R_{\ge0}+\sum_{i=1}^{s}\Sigma U_i^2,
\]
where $P$ is the full cone of nonnegative affine quantum slacks.
Both inclusions must be proved for those actual data.
Take bases $v_{i1},\ldots,v_{id_i}$ and a basis of
\[
 W=\operatorname{span}\{1,\phi_1,\ldots,\phi_8,
                v_{i\alpha}v_{i\beta}:i,\alpha,\beta\}.
\]
For $\ell\in W^*$ impose
\[
 \ell(1)=1,\quad \ell(\phi_\nu)=q_\nu,\qquad
 M_i(\ell)=[\ell(v_{i\alpha}v_{i\beta})]_{\alpha,\beta}
                         \succeq0\quad(1\le i\le s).
\]
The printed product coefficients in the basis of $W$ are exactly
the affine-pencil entries. Evaluation at a covering point proves
the forward inclusion. If a feasible $q$ were outside the compact
convex quantum set, a nonnegative affine slack $f$ with $f(q)<0$
would give
$\ell(\phi^*f)\ge0$ from the SOS identity, a contradiction.
This proves the reverse inclusion without replacing the required
global SOS identity by tests at finitely many real sources.

This construction has total block order $N=\sum_i d_i$ and at
most $\dim W-9$ auxiliary coordinates (the full quantum set has
affine dimension eight). Its entries become evaluated only when
the cover, features and product basis are evaluated. Neither the
local rank bound twelve nor the 377-function local union
supplies the missing global list $U_1,\ldots,U_s$.

\end{document}
